\subsection{Модуль lab1}

\subsubsection{Эксцентриситеты, центр и диаметральные вершины}

\textbf{Вход:} \texttt{Graph<T>}~-- связный граф.

\textbf{Выход:} вектор эксцентриситетов всех вершин; список центральных
вершин (минимальный эксцентриситет); список диаметральных вершин
(максимальный эксцентриситет).

\textbf{Описание:} Для каждой вершины запускается BFS по матрице весов~--
ребро считается существующим, если вес ненулевой. BFS возвращает
расстояния в рёбрах. Эксцентриситет~-- наибольшее расстояние до других вершин.
Центр~-- вершины с минимальным эксцентриситетом, диаметральные~-- с максимальным.

\begin{lstlisting}[style=cstyle, caption={BFS и вычисление эксцентриситета (graph.cpp)}]
template <typename T>
int Graph<T>::eccentricity(int vertex) const {
    vector<int> dist = bfs(vertex);
    const int INF = std::numeric_limits<int>::max();
    int maxDist = 0;
    for (int i = 0; i < numVertices; i++) {
        if (i == vertex || dist[i] == INF) continue;
        maxDist = std::max(maxDist, dist[i]);
    }
    return maxDist;
}
\end{lstlisting}

\begin{lstlisting}[style=cstyle, caption={Поиск центра и диаметральных вершин (graph.cpp)}]
template <typename T>
vector<int> Graph<T>::findCenter() const {
    vector<int> eccs = allEccentricities();
    int radius = *std::min_element(eccs.begin(), eccs.end());
    vector<int> centers;
    for (int i = 0; i < numVertices; i++)
        if (eccs[i] == radius) centers.push_back(i);
    return centers;
}

template <typename T>
vector<int> Graph<T>::findDiametral() const {
    vector<int> eccs = allEccentricities();
    int diameter = 0;
    for (int ecc : eccs)
        if (ecc != std::numeric_limits<int>::max() && ecc > diameter)
            diameter = ecc;
    vector<int> diametral;
    for (int i = 0; i < numVertices; i++)
        if (eccs[i] == diameter) diametral.push_back(i);
    return diametral;
}
\end{lstlisting}

\textbf{Сложность:} $O(V^3)$ по времени (BFS из каждой вершины стоит
$O(V^2)$ при матричном представлении); $O(V)$ по памяти.

\subsubsection{ShimbelSolver\textless T\textgreater: метод Шимбелла}

\textbf{Вход:} \texttt{Graph<T> const\&}, \texttt{int k}~-- длина пути
в рёбрах.

\textbf{Выход:} матрица минимальных и матрица максимальных путей
ровно из $k$ рёбер для всех пар вершин.

\textbf{Описание:} Начальная матрица заполняется по весам графа: если
ребро $(i,j)$ есть~-- ставится $w(i,j)$, иначе~-- INF. Матрица умножается
$k-1$ раз: для каждой пары $(i,j)$ перебираются промежуточные вершины $m$
и выбирается минимальный или максимальный кандидат
$\texttt{mat}[i][m] + w[m][j]$.

\begin{lstlisting}[style=cstyle, caption={Матричное умножение для метода Шимбелла (shimbel.cpp)}]
template <typename T>
void ShimbelSolver<T>::multiplyMatrix(bool findMin) {
    const T INF = getINF();
    const T empty = findMin ? INF : -INF;
    auto weight = graph.getWeightMatrix();
    vector<vector<T>> next(n, vector<T>(n, empty));
    for (int i = 0; i < n; i++)
        for (int j = 0; j < n; j++)
            for (int m = 0; m < n; m++) {
                if (currentMatrix[i][m] == empty
                    || weight[m][j] == T{}) continue;
                T candidate = currentMatrix[i][m] + weight[m][j];
                if (findMin) next[i][j] = std::min(next[i][j], candidate);
                else         next[i][j] = std::max(next[i][j], candidate);
            }
    currentMatrix = next;
}
\end{lstlisting}

\textbf{Сложность:} $O(k \cdot V^3)$ по времени; $O(V^2)$ по памяти.

\subsubsection{Поиск всех маршрутов (DFS с возвратом)}

\textbf{Вход:} \texttt{Graph<T> const\&}, вершина \texttt{from},
вершина \texttt{to}.

\textbf{Выход:} \texttt{vector<vector<int>>}~-- список всех простых
путей из \texttt{from} в \texttt{to}.

\textbf{Описание:} \texttt{findAllPaths} запускает рекурсивный DFS
с возвратом. При достижении целевой вершины текущий
путь сохраняется в результат.
\begin{lstlisting}[style=cstyle, caption={Поиск всех путей DFS (graph.cpp)}]
template <typename T>
void Graph<T>::dfs(int current, int target,
                   vector<bool>& visited,
                   vector<int>& currentPath,
                   vector<vector<int>>& allPaths) const {
    if (current == target) {
        allPaths.push_back(currentPath);
        return;
    }
    for (int next = 0; next < numVertices; next++) {
        if (weightMatrix[current][next] != T{} && !visited[next]) {
            visited[next] = true;
            currentPath.push_back(next);
            dfs(next, target, visited, currentPath, allPaths);
            currentPath.pop_back();
            visited[next] = false;
        }
    }
}
\end{lstlisting}

\textbf{Сложность:} $O(V!)$ по времени в худшем случае;
$O(V)$ по памяти (стек рекурсии и текущий путь).
